% Banco de preguntas de prerrequisitos matemáticos para la defensa.
% Preguntas con respuesta corta, por nivel, con referencia a la tesis.
% Compilar dos veces:  lualatex PreguntasPrerrequisitos.tex
\documentclass[11pt,a4paper]{article}
\usepackage[a4paper,margin=2.3cm,headheight=14pt]{geometry}
\usepackage{amsmath,mathtools}
\usepackage{fontspec}
\usepackage{unicode-math}
\usepackage[spanish,es-noshorthands,es-tabla]{babel}
\setmainfont{TeX Gyre Pagella}
\setmathfont{TeX Gyre Pagella Math}
\setsansfont{DejaVu Sans}[Scale=MatchLowercase]
\usepackage{xcolor}
\usepackage[most]{tcolorbox}
\usepackage{booktabs,tabularx,array}
\usepackage{enumitem}
\usepackage{fancyhdr}
\usepackage[hidelinks]{hyperref}
\usepackage{microtype}
\usepackage{needspace}
\AtBeginDocument{\providecommand\square{\mdlgwhtsquare}}

\definecolor{nb}{HTML}{5E9A62}
\definecolor{ni}{HTML}{3F6FA8}
\definecolor{na}{HTML}{B5523B}
\definecolor{fondonota}{HTML}{F1EEF8}
\definecolor{bordenota}{HTML}{8C7DB8}

\newtcolorbox{nota}[1][Nota]{enhanced, breakable, colback=fondonota,
  colframe=bordenota, boxrule=0.5pt, arc=3pt, colbacktitle=bordenota,
  coltitle=white, fonttitle=\sffamily\bfseries\small, title={#1},
  before skip=8pt, after skip=8pt}

\newcounter{preg}
\newcommand{\nivel}[1]{%
  \ifx#1B\tcbox[on line, boxsep=0pt, left=3pt, right=3pt, top=1pt, bottom=1pt,
    colback=nb!15, colframe=nb, boxrule=0.4pt, arc=2pt]{\sffamily\scriptsize básico}\fi
  \ifx#1I\tcbox[on line, boxsep=0pt, left=3pt, right=3pt, top=1pt, bottom=1pt,
    colback=ni!15, colframe=ni, boxrule=0.4pt, arc=2pt]{\sffamily\scriptsize intermedio}\fi
  \ifx#1A\tcbox[on line, boxsep=0pt, left=3pt, right=3pt, top=1pt, bottom=1pt,
    colback=na!15, colframe=na, boxrule=0.4pt, arc=2pt]{\sffamily\scriptsize avanzado}\fi}
% \Q{nivel}{pregunta}{respuesta}{dónde en la tesis}
\newcommand{\Q}[4]{\par\Needspace{4\baselineskip}\medskip\noindent
  \refstepcounter{preg}\textbf{\thepreg.}\ \nivel{#1}\ \textbf{#2}\par
  \nopagebreak\noindent #3%
  \if\relax\detokenize{#4}\relax\else\ {\small\textcolor{gray}{[#4]}}\fi\par}

\newcommand{\Provnd}{\vdash_{\mathrm{ND}}}
\newcommand{\bred}{\rightarrow_\beta}
\newcommand{\bmred}{\twoheadrightarrow_\beta}
\newcommand{\concat}{\mathbin{+\!\!+}}
\newcommand{\sem}[1]{[\![#1]\!]}

\pagestyle{fancy}
\fancyhf{}
\fancyhead[L]{\small\sffamily Preguntas de prerrequisitos matemáticos}
\fancyhead[R]{\small\sffamily\nouppercase{\leftmark}}
\fancyfoot[C]{\small\thepage}
\renewcommand{\headrulewidth}{0.4pt}
\setlist{itemsep=2pt, topsep=4pt}
\setlength{\parskip}{3pt}
\setlength{\parindent}{0pt}

\title{\textbf{Preguntas de prerrequisitos matemáticos}\\[4pt]
  \large Lo que conviene saber «por si acaso» antes de la defensa}
\author{Material de apoyo para la defensa}
\date{}

\begin{document}
\maketitle
\thispagestyle{empty}

\begin{abstract}
\noindent
Banco de preguntas con respuesta corta sobre las nociones matemáticas que la
tesis da por sabidas o que un tribunal puede usar para comprobar el dominio
del tema: conjuntos, relaciones, inducción, lógica, cálculo~\(\lambda\),
tipos, fundamentos y algoritmos. Cada pregunta lleva un nivel
(\nivel{B}, \nivel{I}, \nivel{A}) y, entre corchetes, dónde aparece la idea en
la tesis. Algunas preguntas tratan temas \emph{cercanos} que la tesis no
desarrolla (Gödel, compacidad, complejidad); están marcadas como tales.
\end{abstract}

\begin{nota}[Cómo usar este banco]
\begin{enumerate}
  \item Tapar las respuestas e intentar contestar en voz alta, en una o dos
    frases.
  \item Si falla alguna pregunta \nivel{B}, repasar esa sección antes de
    seguir: el tribunal da por sabido ese nivel.
  \item Las \nivel{I} son las que se esperan de quien escribió la tesis.
  \item Las \nivel{A} sirven para no quedarse en blanco ante una pregunta
    ambiciosa; basta una respuesta correcta y breve.
\end{enumerate}
\end{nota}

\tableofcontents
\newpage

%=====================================================================
\section{Conjuntos, funciones y listas}
%=====================================================================

\Q{B}{¿Qué significan \(\Gamma\cup\{\varphi\}\), \(\Gamma\subseteq\Delta\) y \(\varnothing\)?}
{Unión: los elementos de \(\Gamma\) más \(\varphi\). Inclusión: todo elemento de
\(\Gamma\) está en \(\Delta\). \(\varnothing\): el conjunto vacío, sin
elementos. En un secuente, \(\Gamma,\varphi\) abrevia \(\Gamma\cup\{\varphi\}\).}
{§11.3.2}

\Q{B}{¿Qué diferencia hay entre conjunto, lista y multiconjunto?}
{Un \emph{conjunto} no tiene orden ni repeticiones. Una \emph{lista} tiene
orden y puede repetir. Un \emph{multiconjunto} no tiene orden pero cuenta
repeticiones. En la tesis: los contextos son conjuntos; los átomos de una
fórmula, una lista; dos listas son permutación una de otra si tienen el mismo
multiconjunto.}
{Caps.\ 11 y 12}

\Q{B}{¿Qué es una función? ¿Qué es una función total y una parcial?}
{Una regla que asigna a cada elemento del dominio un único elemento del
codominio. Total: definida en todo el dominio. Parcial: puede no estar
definida en algunos puntos, por ejemplo porque un cálculo no termina.}
{Prop.\ 12.6}

\Q{B}{¿Qué son una función inyectiva, una sobreyectiva y una biyectiva?}
{Inyectiva: elementos distintos van a imágenes distintas. Sobreyectiva: toda
imagen posible se alcanza. Biyectiva: ambas cosas; tiene inversa. Un
\emph{isomorfismo} es una biyección que respeta la estructura. Por eso
«isomorfismo de Curry--Howard»: pruebas y términos se corresponden
respetando las reglas.}
{Teo.\ 9.22}

\Q{B}{¿Qué significa que un conjunto sea numerable?}
{Que se puede poner en biyección con \(\mathbb N\) (o es finito): sus elementos
se pueden enumerar. El conjunto de variables proposicionales es numerable.}
{§11.2.1}

\Q{B}{¿Qué son la cabeza, la cola, la concatenación y la longitud de una lista?}
{En \(a::r\), \(a\) es la cabeza y \(r\) la cola. \(\ell_1\concat\ell_2\) pega
dos listas. \(|\ell|\) cuenta sus elementos, con repeticiones:
\(|[1,1,2]|=3\).}
{§12.2}

\Q{B}{¿Qué es un predicado?}
{Una propiedad de los elementos: una función que a cada \(x\) le asigna
verdadero o falso. En Quicksort, «\(x\leq p\)» es un predicado y el filtro
conserva los elementos que lo cumplen.}
{Def.\ 12.4}

\Q{I}{¿Qué diferencia hay entre un valor booleano y una proposición?}
{Un booleano es un dato, \(\mathsf V\) o \(\mathsf F\), que se puede calcular.
Una proposición es un enunciado que puede ser demostrable o no, y no siempre
hay un algoritmo que decida cuál. \(\mathsf{decidir}(x\leq p)\) convierte una
proposición decidible en un booleano.}
{§12.2, Nota 11.7}

\Q{I}{¿Qué significa «definido por recursión»?}
{Definir una función dando su valor en los casos básicos y, en los casos
compuestos, en términos de su valor en partes más pequeñas. Ejemplos:
\(\mathrm{FV}\), la sustitución, \(\sem{\cdot}_v\), \(\mathrm{atomos}\) y el filtro.}
{Caps.\ 7, 11 y 12}

%=====================================================================
\section{Relaciones y órdenes}
%=====================================================================

\Q{B}{¿Qué son las propiedades reflexiva, simétrica, transitiva, antisimétrica y total?}
{Reflexiva: \(a\,R\,a\). Simétrica: \(a\,R\,b\Rightarrow b\,R\,a\). Transitiva:
\(a\,R\,b\wedge b\,R\,c\Rightarrow a\,R\,c\). Antisimétrica:
\(a\,R\,b\wedge b\,R\,a\Rightarrow a=b\). Total: \(a\,R\,b\vee b\,R\,a\).}
{§12.2}

\Q{B}{¿Qué es una relación de equivalencia? Dé tres ejemplos de la tesis.}
{Una relación reflexiva, simétrica y transitiva. Ejemplos: \(=_\alpha\)
(mismos términos salvo nombres ligados), \(=_\beta\) (convertibles) y
\(\sim\) (permutación de listas).}
{Def.\ 7.16, Lema 7.32, Def.\ 12.3}

\Q{I}{¿Qué es una clase de equivalencia y qué es un cociente?}
{La clase de \(a\) es el conjunto de todo lo equivalente a \(a\). El cociente
es el conjunto de todas las clases. «Trabajar módulo \(\alpha\)» es
trabajar en el cociente de los términos por \(=_\alpha\). En el asistente,
los cocientes son primitivos, y uno de los tres axiomas de la auditoría
(solidez de los cocientes) dice que elementos equivalentes dan la misma clase.}
{§7.2.4, Nota 10.7}

\Q{B}{¿Qué diferencia hay entre un orden parcial y un orden lineal?}
{Parcial: reflexivo, antisimétrico y transitivo; puede haber elementos
incomparables (\(\subseteq\) entre conjuntos). Lineal o total: además, dos
elementos cualesquiera son comparables (\(\leq\) en \(\mathbb N\)). Quicksort
pide un orden lineal.}
{§12.2}

\Q{B}{¿Cómo se relacionan \(\leq\) y \(<\)?}
{\(a<b\iff a\leq b\wedge a\neq b\). En un orden lineal además
\(\neg(b\leq a)\iff a<b\). Esta última ley garantiza que lo que no va a la
izquierda en Quicksort es mayor que el pivote.}
{(O) en §12.2}

\Q{I}{¿Qué es un preorden? Dé un ejemplo.}
{Una relación reflexiva y transitiva, no necesariamente antisimétrica.
Ejemplo: comparar personas por su edad; dos personas distintas pueden tener
la misma edad.}
{}

\Q{I}{¿Qué es la clausura reflexiva y transitiva de una relación?}
{La menor relación reflexiva y transitiva que la contiene: «cero o más
pasos». \(\bmred\) es la de \(\bred\).}
{Def.\ 7.27}

\Q{I}{¿Qué es una relación bien fundada? Un ejemplo y un contraejemplo.}
{Una relación sin cadenas descendentes infinitas. \((\mathbb N,<)\) es bien
fundada. \((\mathbb Z,<)\) no lo es (\(0>-1>-2>\cdots\)), y tampoco los
racionales no negativos con \(<\) (\(1>\tfrac12>\tfrac14>\cdots\)).}
{§12.4}

\Q{I}{¿Qué significa que una relación sea compatible, o una congruencia?}
{Que se conserva al construir términos más grandes: si \(M\,R\,N\), entonces
\(M\,L\,R\,N\,L\), \(L\,M\,R\,L\,N\) y \(\lambda x.M\,R\,\lambda x.N\). La
reducción \(\beta\) y la \(\alpha\)-equivalencia son compatibles.}
{Defs.\ 7.16, 7.24}

\Q{A}{¿Qué diferencia hay entre la propiedad del diamante, la confluencia y la confluencia local?}
{Diamante: dos pasos que divergen se cierran en \emph{un} paso cada uno.
Confluencia: dos reducciones de cualquier longitud se cierran. Local: dos
pasos se cierran con reducciones de cualquier longitud. La local solo
implica la confluencia si además hay terminación (lema de Newman).}
{§7.2.8}

%=====================================================================
\section{Inducción y recursión}
%=====================================================================

\Q{B}{Enuncie el principio de inducción sobre \(\mathbb N\). ¿Por qué es válido?}
{Si \(P(0)\) y \(P(n)\Rightarrow P(n+1)\) para todo \(n\), entonces
\(P(n)\) para todo \(n\). Es válido porque \(\mathbb N\) es el menor conjunto
que contiene \(0\) y es cerrado bajo sucesor; equivalentemente, todo
subconjunto no vacío de \(\mathbb N\) tiene mínimo.}
{Nota 10.1}

\Q{B}{¿Qué es la inducción fuerte?}
{Suponer \(P(k)\) para \emph{todo} \(k<n\) para probar \(P(n)\). Es
equivalente a la inducción ordinaria. Es la forma que usa Quicksort, sobre la
longitud.}
{§12.4}

\Q{B}{¿Qué son el caso base, el paso inductivo y la hipótesis inductiva?}
{Caso base: los objetos más simples. Paso inductivo: pasar de objetos más
pequeños a uno más grande. Hipótesis inductiva: lo que se supone sobre los
más pequeños durante el paso.}
{Nota 7.5}

\Q{B}{¿Qué es la inducción estructural?}
{Inducción siguiendo la definición de un tipo de objetos: un caso por cada
forma de construirlos (por ejemplo átomo, \(\neg\) y \(\to\)), con hipótesis
inductiva para las partes inmediatas.}
{Nota 7.5, Lema 11.15}

\Q{I}{¿Qué es la inducción sobre derivaciones y en qué se distingue de la estructural?}
{Es la misma idea aplicada a las derivaciones: un caso por cada regla, con
hipótesis inductiva para las premisas. Se usa para probar propiedades de todo
lo \emph{derivable} (debilitamiento, corrección).}
{Lemas 7.57, 11.9; Teo.\ A.3}

\Q{I}{¿Qué es la inducción bien fundada y qué es una medida?}
{Inducción sobre una relación bien fundada: se supone la propiedad para todo
lo «más pequeño». Una medida es una función hacia \(\mathbb N\) que define
«más pequeño». Quicksort usa la longitud de la lista.}
{§12.4}

\Q{I}{¿Qué quiere decir «generalizar la hipótesis» antes de inducir?}
{Dejar cuantificada una variable que cambia de un caso a otro, para que la
hipótesis inductiva pueda usarse con otro valor. En el debilitamiento se
generaliza el contexto \(\Delta\); en la corrección, la valuación \(v\).}
{Lema 11.9, Teo.\ A.3}

\Q{I}{¿Qué quiere decir «reforzar el enunciado»?}
{Demostrar algo más fuerte que lo que se quiere, porque el enunciado original
no pasa la inducción. Ejemplo: el lema de Kalmár es bilateral (verdadero y
falso) aunque la completitud solo use la parte verdadera.}
{Lema 11.15}

\Q{B}{¿Qué es una definición por recursión primitiva? Dé el factorial.}
{Una función con \(g(0)=x\) y \(g(n+1)=f(n,g(n))\). Factorial:
\(0!=1\), \((n+1)!=(n+1)\cdot n!\). El teorema de recursión garantiza que esa
función existe y es única.}
{Nota 10.1}

\Q{I}{¿Qué diferencia hay entre recursión estructural y recursión bien fundada?}
{Estructural: cada llamada es sobre una parte inmediata del argumento, como
la cola; termina automáticamente. Bien fundada: las llamadas son sobre algo
«más pequeño» según una medida; hay que probar que la medida decrece.
Quicksort es del segundo tipo, porque \(L\) y \(R\) no son la cola.}
{§12.4}

\Q{A}{¿Por qué una «definición» recursiva que no termina es un problema en lógica?}
{Porque sus ecuaciones pueden ser contradictorias. Si se admitiera
\(g(n)=g(n)+1\), restando se obtendría \(0=1\). Por eso los asistentes exigen
probar la terminación.}
{Prop.\ 12.6}

\Q{I}{¿Qué es un invariante?}
{Una propiedad que se conserva en cada paso de un proceso. En Quicksort,
«la partición no pierde ni duplica elementos» y «todo lo de la izquierda es
\(\leq p\)» son invariantes del paso de partición.}
{Conclusión del Cap.\ 12}

%=====================================================================
\section{Lógica proposicional: sintaxis y semántica}
%=====================================================================

\Q{B}{¿Qué son un átomo, un conectivo y una fórmula?}
{Átomo: una variable proposicional (\(p\)). Conectivo: un símbolo que
combina fórmulas (\(\neg,\to,\wedge,\vee\)). Fórmula: una expresión bien
formada construida a partir de átomos con conectivos.}
{Def.\ 11.3}

\Q{B}{¿Cómo se lee \(\varphi::=p\mid\neg\varphi\mid(\varphi\to\psi)\)?}
{Es una gramática (notación BNF): una fórmula es un átomo, \emph{o} la
negación de una fórmula, \emph{o} la implicación entre dos fórmulas. Cada
\(\mid\) separa una alternativa.}
{Def.\ 11.3}

\Q{B}{¿Qué es una valuación? ¿Cuántas filas tiene la tabla de una fórmula con \(n\) átomos?}
{Una asignación de \(\mathsf V\) o \(\mathsf F\) a cada átomo; cada fila de la
tabla es una valuación. Con \(n\) átomos hay \(2^n\) filas.}
{Def.\ 11.4}

\Q{B}{¿Por qué «falso implica cualquier cosa» es verdadero?}
{En la implicación material, \(\varphi\to\psi\) solo es falsa cuando
\(\varphi\) es verdadera y \(\psi\) falsa. Una promesa «si llueve, llevo
paraguas» no se incumple si no llueve.}
{§9.2}

\Q{B}{¿Qué son una tautología, una contradicción y una fórmula satisfacible?}
{Tautología: verdadera en todas las filas. Contradicción: falsa en todas.
Satisfacible: verdadera en alguna. Toda tautología es satisfacible; una
contradicción no.}
{Def.\ 9.12}

\Q{B}{¿Qué significa \(\Gamma\models\varphi\)? ¿Cómo se relaciona con las tautologías?}
{Toda valuación que hace verdaderas las fórmulas de \(\Gamma\) hace
verdadera \(\varphi\). Si \(\Gamma=\{\gamma_1,\dots,\gamma_n\}\), entonces
\(\Gamma\models\varphi\) equivale a que
\(\gamma_1\to(\cdots\to(\gamma_n\to\varphi))\) sea tautología.}
{Def.\ 9.11, Teo.\ 9.14}

\Q{B}{¿Qué es un literal?}
{Un átomo o la negación de un átomo: \(p\) o \(\neg p\). En el argumento de
Kalmár, el literal de \(p\) en una fila es lo que esa fila dice de \(p\).}
{Def.\ 11.13}

\Q{I}{¿Qué es una función de verdad? ¿Cuántas hay de \(n\) argumentos?}
{Una función \(\{\mathsf V,\mathsf F\}^n\to\{\mathsf V,\mathsf F\}\), es decir,
una tabla de verdad. Hay \(2^{2^n}\): \(2^n\) filas y dos valores posibles en
cada una. Para \(n=2\), hay 16.}
{§11.2.2}

\Q{I}{¿Qué significa «base» de conectivos?}
{El conjunto de conectivos que se toman como \emph{primitivos}; los demás se
definen como abreviaturas. La tesis usa la base \(\{\neg,\to\}\).}
{§11.2.2}

\Q{I}{¿Qué significa que una base sea «funcionalmente completa»?}
{Que toda función de verdad, de cualquier número de argumentos, puede
escribirse como una fórmula que solo usa esos conectivos. Es una propiedad de
\emph{expresividad}, no de demostrabilidad.}
{§11.2.2}

\Q{I}{¿Por qué \(\{\neg,\to\}\) es funcionalmente completa?}
{Toda tabla se escribe en forma normal disyuntiva (una conjunción de
literales por fila verdadera), que solo usa \(\neg,\wedge,\vee\). A su vez,
\(A\vee B\equiv\neg A\to B\) y \(A\wedge B\equiv\neg(A\to\neg B)\).}
{§11.2.2; Mendelson, cor.\ 1.6}

\Q{I}{¿Por qué \(\{\wedge,\vee\}\) no es funcionalmente completa?}
{Porque cualquier fórmula hecha solo con \(\wedge\) y \(\vee\) es verdadera
cuando todos sus átomos lo son. Entonces no puede expresar \(\neg p\), que
es falsa cuando \(p\) es verdadera. El mismo argumento vale para
\(\{\to\}\) y para \(\{\wedge,\vee,\to\}\).}
{}

\Q{I}{¿Qué son la forma normal disyuntiva y la conjuntiva?}
{Disyuntiva: una disyunción de conjunciones de literales. Conjuntiva: una
conjunción de disyunciones de literales. Toda fórmula es equivalente a una
de cada tipo.}
{§11.2.2}

\Q{A}{¿Existe un único conectivo funcionalmente completo?}
{Sí: \(\mathsf{NAND}\) (no-y) y \(\mathsf{NOR}\) (no-o), llamados conectivos de
Sheffer. Por ejemplo, \(\neg p\equiv p\ \mathsf{NAND}\ p\).}
{}

\Q{I}{¿Qué diferencia hay entre \(\equiv\) (equivalencia lógica), \(\dashv\vdash\) y la igualdad sintáctica?}
{\(\varphi\equiv\psi\): mismo valor en todas las filas (semántica).
\(\varphi\dashv\vdash\psi\): cada una se deriva de la otra (sintaxis). Igualdad
sintáctica: son la misma expresión. Por adecuación, las dos primeras
coinciden en lógica clásica.}
{Def.\ 9.8, §11.2.2}

\Q{I}{¿Por qué \(\neg(p_0\to p_0)\) sirve como \(\bot\)?}
{\(p_0\to p_0\) es verdadera en todas las filas, así que su negación es falsa
en todas, cualquiera que sea el valor de \(p_0\).}
{Def.\ 11.8, Lema A.1}

%=====================================================================
\section{Sistemas de demostración}
%=====================================================================

\Q{B}{¿Qué es una regla de inferencia?}
{Un esquema «de estas premisas se concluye esto», escrito como fracción:
premisas arriba y conclusión abajo.}
{§9.1}

\Q{B}{¿Qué es una derivación, o prueba formal?}
{Un árbol de aplicaciones de reglas cuya raíz es la conclusión y cuyas hojas
son hipótesis o axiomas.}
{§9.1, §11.3}

\Q{B}{¿Qué es un secuente y qué es un teorema del cálculo?}
{\(\Gamma\vdash\varphi\): «\(\varphi\) se deriva de las hipótesis \(\Gamma\)».
Un teorema es una fórmula derivable sin hipótesis: \(\vdash\varphi\).}
{Defs.\ 9.3, 9.4}

\Q{B}{¿Qué diferencia hay entre reglas de introducción y de eliminación?}
{Introducción: cómo \emph{probar} una fórmula con ese conectivo (\(\to\)I).
Eliminación: cómo \emph{usarla} (\(\to\)E, el modus ponens).}
{§9.1}

\Q{B}{¿Qué es una suposición y qué significa descargarla?}
{Una hipótesis temporal que se abre dentro de la prueba («supongamos
\(\varphi\)»). Descargarla es cerrarla al aplicar una regla como \(\to\)I, que
la incorpora a la conclusión (\(\varphi\to\psi\)). Después ya no está
disponible.}
{Nota 9.5}

\Q{B}{¿Qué es el modus ponens?}
{De \(\varphi\) y \(\varphi\to\psi\) se concluye \(\psi\). Es \(\to\)E.}
{§9.1}

\Q{I}{¿En qué se diferencian \(\neg\)I y RAA?}
{\(\neg\)I: de \(\Gamma,A\vdash\bot\) se concluye \(\neg A\); es
intuicionista. RAA: de \(\Gamma,\neg A\vdash\bot\) se concluye \(A\); es
clásica. La forma es parecida; la diferencia es qué se supone y qué se
concluye.}
{Nota 7.60, §11.3.3}

\Q{I}{¿Qué es el principio de explosión?}
{De una contradicción se sigue cualquier cosa: \(\bot\vdash\varphi\). Es la
regla \(\bot\)E.}
{§9.1}

\Q{I}{¿Qué diferencia hay entre una regla primitiva, una derivada y una admisible?}
{Primitiva: forma parte del sistema. Derivada: abrevia una derivación hecha
con las primitivas (el análisis por casos, Lema 11.11). Admisible: añadirla
no cambia lo que se deriva; se prueba por inducción sobre las derivaciones
(el debilitamiento, Lema 11.9).}
{Caps.\ 9 y 11}

\Q{I}{¿Qué diferencia hay entre un sistema de Hilbert y la deducción natural?}
{Hilbert: muchos axiomas y una sola regla (modus ponens), sin suposiciones
temporales. Deducción natural: pocas reglas por conectivo, con suposiciones
que se abren y se descargan. Mendelson usa Hilbert; la tesis usa deducción
natural.}
{§11.1}

\Q{I}{¿Qué dice el teorema de la deducción?}
{\(\Gamma,\varphi\vdash\psi\) si y solo si \(\Gamma\vdash\varphi\to\psi\). En los
sistemas de Hilbert es un teorema; en la deducción natural son directamente
las reglas \(\to\)I y \(\to\)E.}
{}

\Q{I}{¿Qué son las reglas estructurales?}
{Las que manipulan el contexto: debilitamiento (añadir hipótesis),
intercambio (reordenarlas) y contracción (eliminar repeticiones). Con
contextos como conjuntos, las dos últimas vienen gratis; el debilitamiento
se demuestra.}
{§11.3.2}

\Q{A}{¿Qué es un desvío en una prueba y qué es normalizarla?}
{Un desvío es introducir un conectivo e inmediatamente eliminarlo.
Normalizar es eliminar desvíos. Bajo Curry--Howard es exactamente la
reducción \(\beta\).}
{§9.5.2}

%=====================================================================
\section{Metateoremas: corrección, completitud y compañía}
%=====================================================================

\Q{B}{¿Qué es un metateorema?}
{Un teorema \emph{sobre} un sistema formal, no dentro de él. Ejemplos:
corrección, completitud y debilitamiento.}
{Nota 11.1}

\Q{B}{¿Qué dicen la corrección, la completitud y la adecuación?}
{Corrección: si \(\vdash\varphi\), entonces \(\models\varphi\). Completitud: si
\(\models\varphi\), entonces \(\vdash\varphi\). Adecuación: las dos juntas.}
{Teos.\ A.3, 11.18, A.6}

\Q{B}{¿Qué es la consistencia y qué relación tiene con la corrección?}
{Un cálculo es consistente si no deriva \(\bot\). Se sigue de la corrección:
\(\bot\) no es verdadera en ninguna fila, luego no es derivable.}
{Cor.\ A.5}

\Q{I}{¿Es decidible la lógica proposicional?}
{Sí: para saber si \(\varphi\) es tautología basta revisar sus \(2^n\) filas.
No se conoce un método eficiente en general (decidir tautologías es
coNP-completo); esto es contexto, no parte de la tesis.}
{}

\Q{I}{¿Qué diferencia hay entre completitud débil y fuerte?}
{Débil: toda tautología es derivable (\(\Gamma=\varnothing\)); es la que
demuestra la tesis. Fuerte: \(\Gamma\models\varphi\Rightarrow\Gamma\vdash\varphi\)
para todo \(\Gamma\). Para \(\Gamma\) finito se obtiene de la débil con
\(\to\)E; para \(\Gamma\) infinito hace falta la compacidad.}
{Teo.\ 11.18}

\Q{A}{¿Qué es la compacidad? (cercano a la tesis)}
{Un conjunto de fórmulas es satisfacible si y solo si cada subconjunto
finito lo es. No se usa en la tesis porque se trabaja con contextos finitos.}
{}

\Q{A}{¿Qué tiene que ver esto con los teoremas de Gödel?}
{Hay que distinguir tres resultados. Post (1921): completitud de la lógica
\emph{proposicional}, el tema de la tesis. Gödel (1929): completitud de la
lógica de \emph{primer orden}. Gödel (1931): \emph{incompletitud} de la
aritmética, que es otra cosa (ninguna teoría consistente, recursivamente
axiomatizable y que contenga la aritmética demuestra todas las verdades
aritméticas). La tesis no trata los dos últimos.}
{§11.1}

\Q{I}{¿Qué diferencia hay entre el método de Kalmár y el de Lindenbaum--Henkin?}
{Kalmár recorre la tabla de verdad, que es finita, y construye una derivación
por fila. Lindenbaum--Henkin extiende un conjunto consistente a uno maximal y
construye un modelo con él; es el método para primer orden. La tesis usa
Kalmár.}
{§11.5}

\Q{A}{¿Por qué la deducción natural intuicionista no es completa respecto de las tablas de verdad?}
{Porque \(p\vee\neg p\) es tautología y no es derivable sin una regla
clásica. Para la lógica intuicionista, la semántica adecuada es la de
Kripke.}
{Nota 9.18}

%=====================================================================
\section{Lógica clásica e intuicionista}
%=====================================================================

\Q{B}{¿Qué son LEM, la eliminación de la doble negación y RAA?}
{LEM: \(\varphi\vee\neg\varphi\). Eliminación de la doble negación:
\(\neg\neg\varphi\vdash\varphi\). RAA: de \(\neg\varphi\vdash\bot\), concluir
\(\varphi\). Añadir cualquiera de las tres al núcleo intuicionista da la
lógica clásica.}
{Nota 9.7}

\Q{B}{¿Qué es la interpretación BHK?}
{Brouwer--Heyting--Kolmogorov: una prueba es una construcción. Una prueba de
\(A\wedge B\) es un par de pruebas; de \(A\vee B\), una prueba de una de las
dos indicando cuál; de \(A\to B\), un método que transforma pruebas de \(A\)
en pruebas de \(B\); de \(\bot\), no hay ninguna.}
{Def.\ 9.19}

\Q{I}{¿Por qué \(p\vee\neg p\) no es válido intuicionistamente?}
{Bajo BHK exigiría un método que, para cualquier \(p\), decida y dé una
prueba de \(p\) o de \(\neg p\). Tal método no existe en general.}
{Ej.\ 9.20}

\Q{I}{¿Cuál de estas es intuicionista: \(p\to\neg\neg p\) o \(\neg\neg p\to p\)?}
{La primera: dada una prueba de \(p\) y un método de \(p\to\bot\), se aplica
el método. La segunda requiere lógica clásica.}
{Ej.\ 9.20, Lema 11.10}

\Q{I}{¿Qué significa que una prueba sea constructiva?}
{Que, cuando demuestra que algo existe, da un testigo o un método para
obtenerlo, y no usa LEM ni sus equivalentes.}
{§9.4}

\Q{I}{¿Es clásica la lógica del asistente usado en la tesis?}
{Su núcleo es constructivo. El razonamiento clásico se obtiene del axioma de
elección: junto con la extensionalidad, permite derivar el tercero excluido.
Por eso la auditoría distingue entre los teoremas que usan elección (la
completitud) y los que no (Quicksort).}
{§11.6, §12.10}

\Q{A}{¿Qué es la semántica de Kripke? (cercano a la tesis)}
{Una semántica para la lógica intuicionista con «estados de conocimiento»
que solo crecen. Una fórmula puede no ser verdadera todavía sin ser falsa;
por eso \(p\vee\neg p\) puede fallar. La tesis la menciona pero no la
formaliza.}
{Nota 9.18}

%=====================================================================
\section{Cálculo lambda}
%=====================================================================

\Q{B}{¿Cuáles son los tres tipos de términos \(\lambda\)?}
{Variables (\(x\)), aplicaciones (\(MN\): aplicar \(M\) a \(N\)) y abstracciones
(\(\lambda x.M\): la función que a \(x\) le asigna \(M\)).}
{Def.\ 7.1}

\Q{B}{¿Qué son una variable libre y una ligada?}
{Ligada: está bajo un \(\lambda\) que la nombra, como \(x\) en \(\lambda x.x\,y\).
Libre: no lo está, como \(y\). \(\mathrm{FV}(\lambda x.x\,y)=\{y\}\).}
{Def.\ 7.10}

\Q{B}{¿Qué es la \(\alpha\)-equivalencia?}
{Igualdad salvo el nombre de las variables ligadas:
\(\lambda x.x=_\alpha\lambda y.y\). No se permite renombrar si se captura una
variable libre.}
{Defs.\ 7.15, 7.16}

\Q{B}{¿Qué es la sustitución y qué es la captura?}
{\(M[x:=N]\) reemplaza las \(x\) libres de \(M\) por \(N\). Hay captura si una
variable libre de \(N\) queda ligada al entrar bajo un \(\lambda\). Se evita
renombrando primero la variable ligada.}
{Defs.\ 7.17, 7.20}

\Q{B}{¿Qué es un paso \(\beta\), un redex y un contractum?}
{\((\lambda x.M)N\bred M[x:=N]\). El redex es \((\lambda x.M)N\); el contractum,
\(M[x:=N]\). Es «aplicar la función».}
{Defs.\ 7.24, 7.25}

\Q{B}{¿Qué es currificar?}
{Representar una función de varios argumentos como una función que devuelve
funciones: \(\lambda x\,y.M=\lambda x.\lambda y.M\). Por eso
\(\alpha\to\beta\to\gamma\) se lee \(\alpha\to(\beta\to\gamma)\).}
{Notas 7.9, 7.46}

\Q{I}{¿Qué es una forma normal y qué es \(\Omega\)?}
{Forma normal: un término sin redexes. \(\Omega=(\lambda x.xx)(\lambda x.xx)\)
se reduce a sí mismo para siempre, así que no tiene forma normal.}
{Def.\ 7.33, Ej.\ 7.36}

\Q{I}{¿Qué diferencia hay entre normalización débil y fuerte?}
{Débil: \emph{alguna} ruta de reducción termina. Fuerte: \emph{todas}
terminan. \((\lambda y.z)\,\Omega\) es débilmente pero no fuertemente
normalizable.}
{Def.\ 7.35}

\Q{I}{¿Qué dice Church--Rosser y qué consecuencia tiene?}
{Si \(M\) se reduce a \(N_1\) y a \(N_2\), ambos se reducen a un \(L\) común.
Consecuencia: la forma normal, si existe, es única. No dice que exista.}
{Teo.\ 7.37, Cor.\ 7.44}

\Q{I}{¿Qué es la convención de Barendregt?}
{Escribir los términos de modo que las variables ligadas sean distintas
entre sí y distintas de las libres. Así la sustitución no necesita
renombrar.}
{§7.2.6}

\Q{I}{¿Qué es la regla \(\eta\)?}
{\(\lambda x.M\,x\) equivale a \(M\) si \(x\) no está libre en \(M\): dos funciones
que dan lo mismo en todo argumento son iguales. El Ejemplo 7.53 tipa la
\(\eta\)-expansión \(\lambda z.y\,z\) de \(y\).}
{Ej.\ 7.53}

\Q{A}{¿Qué es un combinador de punto fijo?}
{Un término \(Y\) con \(Y\,f=_\beta f\,(Y\,f)\), por ejemplo
\(Y=\lambda f.(\lambda x.f(xx))(\lambda x.f(xx))\). Permite la recursión en el
cálculo no tipado. No es tipable en \(\lambda\!\to\), porque allí todo término
tipado termina.}
{}

\Q{A}{¿Qué potencia computacional tiene el cálculo \(\lambda\) no tipado?}
{Es Turing-completo: calcula exactamente las funciones computables. Church
lo usó en 1936 para mostrar que hay problemas indecidibles.}
{§7.2.1}

%=====================================================================
\section{Tipos y Curry--Howard}
%=====================================================================

\Q{B}{¿Qué es un tipo simple?}
{Una variable de tipo (\(\alpha\)) o un tipo flecha (\(\sigma\to\tau\), funciones
de \(\sigma\) en \(\tau\)).}
{Def.\ 7.45}

\Q{B}{¿Qué es un juicio de tipado?}
{\(\Gamma\vdash M:\sigma\): en el contexto \(\Gamma\) (declaraciones
\(x:\sigma\)), el término \(M\) tiene tipo \(\sigma\).}
{Def.\ 7.51}

\Q{I}{¿Qué diferencia hay entre tipado à la Church y à la Curry?}
{Church: los tipos se escriben en el término (\(\lambda x{:}\sigma.M\)).
Curry: el término no lleva tipos y se le asignan después. La tesis usa
principalmente el estilo de Church.}
{§7.3.2}

\Q{I}{¿Qué son los problemas de tipabilidad, chequeo y habitación?}
{Tipabilidad: ¿tiene \(M\) algún tipo? Chequeo: ¿tiene \(M\) este tipo?
Habitación: ¿hay algún \(M\) de este tipo? En \(\lambda\!\to\) los tres son
decidibles; en \(\lambda P\) y \(\lambda C\) la habitación no lo es.}
{§7.3.4, Teo.\ 7.59, §8.4.3}

\Q{I}{¿Por qué \(x\,x\) no es tipable?}
{\(x\) tendría que tener tipo \(\sigma\to\tau\) y tipo \(\sigma\) a la vez, así que
\(\sigma=\sigma\to\tau\), lo cual es imposible por tamaño.}
{Ej.\ 7.49}

\Q{B}{¿Qué dice la correspondencia de Curry--Howard?}
{Proposiciones son tipos y pruebas son programas. \(\to\) es el tipo función,
\(\wedge\) el producto, \(\vee\) la suma. Probar es construir un término;
simplificar una prueba es evaluar el programa.}
{§9.5}

\Q{I}{¿Qué significa «un tipo está habitado»?}
{Que existe un término de ese tipo. Bajo Curry--Howard: que la proposición
correspondiente es demostrable.}
{Def.\ 7.61}

\Q{I}{¿Qué son la preservación del tipo, la unicidad de tipos y la normalización fuerte?}
{Preservación: si \(M:A\) y \(M\bred N\), entonces \(N:A\). Unicidad: un término
tiene un solo tipo salvo \(=_\beta\). Normalización fuerte: todo término tipado
termina. Valen en todo el cubo (citado).}
{Teo.\ 8.27}

\Q{I}{¿Qué es el polimorfismo y qué es \(\Pi\)?}
{Polimorfismo: un término que funciona para todo tipo, como la identidad
\(\Lambda\alpha.\lambda x{:}\alpha.x:\Pi\alpha{:}*.\alpha\to\alpha\). \(\Pi\) es el
tipo de las funciones cuyo tipo de salida depende del argumento.}
{§8.1}

\Q{I}{¿Qué son un kind y un sort?}
{Kind: el tipo de un tipo (\(\mathsf{nat}:*\), \(\mathsf{List}:*\to*\)). Sorts:
las clasificaciones del nivel superior, \(*\) y \(\square\), con \(*:\square\).}
{§8.2}

\Q{I}{¿Qué es un tipo dependiente?}
{Un tipo que depende de un término, como \(\mathsf{Vec}\ n\) (vectores de
longitud \(n\)) o \(\mathsf{EsPar}\ n\). Permite expresar \(\forall\) como \(\Pi\).}
{§8.3}

\Q{A}{¿Cuáles son los tres ejes del cubo \(\lambda\)?}
{Términos que dependen de tipos (\(\lambda2\), polimorfismo); tipos que dependen
de tipos (\(\lambda\underline{\omega}\), constructores); tipos que dependen de
términos (\(\lambda P\), lógica de predicados). \(\lambda C\) tiene los tres.}
{Ej.\ 8.26}

%=====================================================================
\section{CIC y fundamentos}
%=====================================================================

\Q{I}{¿Qué añade el CIC al CoC?}
{Tipos inductivos primitivos: se declaran con sus constructores, y el sistema
genera su eliminador (recursión e inducción) y su regla de cómputo \(\iota\).}
{§10.3}

\Q{I}{¿Qué son un constructor, un eliminador y la reducción \(\iota\)?}
{Constructor: la forma de construir un elemento (\(0\), \(S\)). Eliminador: la
forma de usarlo, que da la recursión y la inducción. Reducción \(\iota\): el
eliminador aplicado a un constructor se calcula con el caso de ese
constructor.}
{§10.3}

\Q{I}{¿Qué es la codificación de Church y por qué no basta?}
{Representar \(n\) como «aplicar \(f\) \(n\) veces». Da iteración, pero de ella no
se deriva el principio de inducción (resultado de Geuvers, citado), y el
predecesor es costoso.}
{§10.2.3}

\Q{A}{¿Qué significa «impredicativo»?}
{Que una definición cuantifica sobre una colección que incluye a lo que se
está definiendo. Por ejemplo, \(\forall A{:}\mathrm{Prop}.\,A\) es una
proposición definida cuantificando sobre todas las proposiciones.}
{§10.2}

\Q{A}{¿Por qué hay una jerarquía de universos y no \(\mathrm{Type}:\mathrm{Type}\)?}
{Porque \(\mathrm{Type}:\mathrm{Type}\) hace inconsistente el sistema (paradoja
de Girard, análoga a la de Russell). Se usa
\(\mathrm{Type}_0:\mathrm{Type}_1:\cdots\).}
{§8.2.1}

\Q{A}{¿Por qué se exige positividad estricta en los tipos inductivos?}
{Porque un constructor como \(c:(X\to\bot)\to X\) permitiría construir un
término diagonal de tipo \(\bot\) (como la paradoja de Russell), y la lógica
sería inconsistente.}
{§10.3.2}

\Q{I}{¿Qué es la irrelevancia de pruebas?}
{Dos pruebas de la misma proposición se consideran iguales: importa
\emph{que} haya prueba, no \emph{cuál}.}
{Nota 10.7}

\Q{I}{¿Qué dicen los tres axiomas que reporta la auditoría?}
{\emph{Extensionalidad proposicional}: dos proposiciones equivalentes son
iguales. \emph{Elección}: de la existencia se puede obtener un testigo; da el
razonamiento clásico. \emph{Solidez de los cocientes}: elementos equivalentes
tienen la misma clase; de ella se deriva la extensionalidad de funciones.}
{§11.6}

\Q{I}{¿Qué son un asistente de pruebas y su núcleo?}
{Un programa que comprueba demostraciones formales. El núcleo es la parte
pequeña que verifica los tipos; solo hay que confiar en él y en los axiomas,
no en la automatización que genera las pruebas (criterio de de Bruijn).}
{Cap.\ 10}

%=====================================================================
\section{Algoritmos}
%=====================================================================

\Q{B}{¿Qué exige la especificación de un algoritmo de ordenamiento?}
{Que la salida sea una permutación de la entrada y que esté en orden no
decreciente. Ninguna de las dos condiciones basta sola.}
{Def.\ 12.1}

\Q{B}{¿Qué es una permutación y qué es la multiplicidad?}
{Multiplicidad: cuántas veces aparece un elemento. Dos listas son
permutación una de otra si cada elemento tiene la misma multiplicidad en
ambas.}
{Def.\ 12.3}

\Q{B}{¿Cómo funciona Quicksort y qué es el pivote?}
{Se toma el primer elemento como pivote, se separa el resto en los \(\leq\)
pivote y los \(>\) pivote, se ordena cada parte recursivamente y se pegan con
el pivote en medio. Es un ejemplo de divide y vencerás.}
{Def.\ 12.4}

\Q{I}{¿Qué diferencia hay entre corrección parcial y total?}
{Parcial: \emph{si} el algoritmo termina, el resultado es correcto. Total:
además termina siempre. La tesis prueba la total (terminación + permutación +
orden).}
{Cap.\ 12}

\Q{I}{¿Cuál es la complejidad de Quicksort? (fuera del alcance)}
{\(\Theta(n\log n)\) en promedio y \(\Theta(n^2)\) en el peor caso, por ejemplo con
una lista ya ordenada y el primer elemento como pivote. La tesis no la
analiza.}
{§12.1}

\Q{I}{¿Qué es una implementación in situ?}
{Una que reordena dentro del mismo arreglo, con memoria auxiliar constante
(la partición de Hoare). La versión de la tesis es funcional: construye
listas nuevas.}
{§12.1}

\Q{I}{¿Qué es el árbol de recursión?}
{Un árbol con un nodo por llamada y un hijo por cada llamada recursiva. En
Quicksort, leerlo en orden simétrico da la salida ordenada.}
{§12.3.2}

\Q{A}{¿Es estable el Quicksort de la tesis? (observación, no está en la tesis)}
{No necesariamente. Las copias del pivote que están en la cola quedan
\emph{antes} que el pivote: si \(2_a\) y \(2_b\) son iguales,
\([2_a,1,2_b]\mapsto[1,2_b,2_a]\). La especificación no pide estabilidad,
así que el teorema no se ve afectado.}
{}

%=====================================================================
\section{Verificación formal}
%=====================================================================

\Q{B}{¿Qué significa que un teorema esté verificado formalmente?}
{Que su prueba fue comprobada paso a paso por el núcleo de un asistente, a
partir de definiciones formales y de un conjunto explícito de axiomas.}
{Caps.\ 11 y 12}

\Q{I}{¿Qué comprueba una auditoría de axiomas?}
{Lista todos los axiomas de los que depende un teorema, directa o
indirectamente. Si aparece una prueba pendiente (un hueco), se ve. En la
tesis solo aparecen los tres axiomas estándar.}
{§11.6, §12.10}

\Q{I}{¿Qué \emph{no} garantiza una verificación formal?}
{Que la definición formal capture lo que se quería decir. Si la
especificación está mal escrita, se demuestra con rigor algo equivocado. Por
eso hay que defender las definiciones (\(\mathsf{ND}\), \(\mathsf{Ordenada}\), la
permutación).}
{}

\Q{A}{¿Por qué confiar en el núcleo?}
{Porque es pequeño, está muy probado y su lógica está descrita con
precisión. La automatización puede tener errores, pero cualquier prueba que
produzca pasa por el núcleo. Es el criterio de de Bruijn.}
{Nota 10.7}

\end{document}
